\documentclass[11pt]{article}
\usepackage[a4paper,margin=25mm]{geometry}
\usepackage{amsmath,amssymb,amsthm}
\usepackage{lmodern}
\usepackage{fontspec}
\setmainfont{lmroman10-regular.otf}[
 BoldFont=lmroman10-bold.otf,
 ItalicFont=lmroman10-italic.otf,
 BoldItalicFont=lmroman10-bolditalic.otf]
\usepackage{tikz-cd}
\usepackage{xurl}
\usepackage[hidelinks]{hyperref}
\newcommand{\rat}{\dashrightarrow}
\newcommand{\Pic}{\operatorname{Pic}}
\newcommand{\NS}{\operatorname{NS}}
\newcommand{\exc}{\mathfrak{e}}
\newcommand{\sourceheading}[1]{\par\medskip\noindent\textbf{#1.} --- }
\newcommand{\sourcestatement}[2]{\par\medskip\noindent\textbf{#1.} --- \textit{#2}\par}
\setlength{\emergencystretch}{2em}
\begin{document}
\begin{center}
 {\Large 3198: Universal minimal blowup stabilisation}\par\medskip
 {\large Richard A. P. Birkett}\par
 \emph{On the Stabilisation of Rational Surface Maps}
\end{center}

\section*{Editorial provenance and scope}
The author's Theorem 1.4 is reproduced below with its complete authored
proof, the complete governing proof of Proposition 1.2, the author's
algorithm, and the necessary surface-map notation and local-stability
context. Supporting statements are not additional corpus items. The
selected conclusion assumes that a stabilising birational morphism exists;
it does not assert existence for every dominant rational map.

\noindent\textbf{Source.} \emph{Annales de la Faculté des sciences de Toulouse},
series 6, volume 34, no. 1 (2025), pp. 47--74.
DOI: \href{https://doi.org/10.5802/afst.1805}{10.5802/afst.1805}.
Official unchanged original:
\url{https://afst.centre-mersenne.org/item/10.5802/afst.1805.pdf}.
The original's individual cover states Creative Commons Attribution 4.0
(CC BY 4.0), copyright the author, 2025.

\noindent\textbf{Source location.} Original physical pages 2--4 and 7--12
(printed pp. 47--49 and 52--57), with bibliography on physical page 29
(printed p. 74). The substantive proof begins on physical page 10 and
finishes on physical page 12, before the subsequent eventual-stability
analogue. The original PDF is identified by the official URL above; its SHA-256 is
\texttt{3e7d00b2cd39095b4dbcf26904dde6207e32e406e4d7d1b8f72865e88b161719}.

\noindent\textbf{Difficulty: H1.} The authored proof forces any stabilising
birational morphism to blow up every point of a minimal destabilising orbit,
using locally stable compositions and total transforms, then turns that
factorisation into bounded termination and uniqueness of the dynamical
model. This combines surface birational geometry with orbit-level algebraic
stability beyond routine graduate qualifying material.

\noindent\textbf{Transcription note.} The following mathematical text is
the author's human-authored text, faithfully transcribed from the indicated
PDF pages; it is not a generated replacement proof. Original statement
numbers, notation, inequalities and three commutative diagrams are retained
in editable TeX. The introduction's algorithm begins with $f_0$, whereas
the proof begins with $(f_1,X_1)$ and retains the original displayed products
and indices. The author's occasional abbreviated point notation, including
$g^{-1}(p_k)$ after defining $\widehat p_k$, is retained. No mathematical
correction is made. Cross-references to other, unselected examples and
remarks retain their original numbers. Abstracts, acknowledgements and
unrelated later primary results are not part of this selected unit.

\clearpage
\section*{Author's text: 1. Introduction}
% SOURCE physical 2, printed 47: beginning of Introduction.
Let $f:X\rat X$ be a rational map on a smooth projective surface over an
algebraically closed field. Studying $f$ as a dynamical system is complicated
by the fact that $f$ need not be continuously defined on all points of $X$.
For example a rational map $f$ induces a natural pullback operator on curves
$f^*:\Pic(X)\to\Pic(X)$, but this operator may not iterate well. We say
$f$ is algebraically stable iff $\forall n\in\mathbb N\ (f^*)^n=(f^n)^*$.
This is equivalent to the geometric condition that $f$ has no destabilising
orbits [13] [7, 1.14]; i.e. an orbit of (closed) points
$p,f(p),\ldots,f^{n-1}(p)$ in $X$, for which $f^{-1}(p)$ and
$f(f^{n-1}(p))$ are curves. It is natural to hope that blowing up the points
in such an orbit will improve the situation. The main theme of this paper
is to discuss the extent to which this actually works.

% SOURCE physical 3, printed 48.
Assume for the rest of the paper that all surfaces are projective varieties
over an algebraically closed field, and unless explicitly stated, also
smooth. We write $\phi:X\rat Y$ to indicate that $f$ is a rational map
between surfaces, and we use a solid arrow $f:X\to Y$ when $f$ is a morphism.

\sourcestatement{Definition 1.1}{We write $\phi:(g,Y)\rat(f,X)$ to indicate
that $\phi:Y\rat X$ is a birational map conjugating $f:X\rat X$ to
$g=\phi^{-1}\circ f\circ\phi:Y\rat Y$. When $g:Y\rat Y$ is algebraically
stable, we say that $\phi$ stabilises $f$.}

Diller--Favre [7] proved that for a birational map $(f,X)$ there is always
a birational morphism $\pi:(g,Y)\to(f,X)$ which stabilises $f$. Not all
(non-invertible) rational maps can be stabilised, however. Favre [9] showed
for example that many monomial maps on $\mathbb P^2$ cannot be stabilised by
any birational conjugacy. See also [18]; for discussion on monomial maps in
higher dimensions see [20, 21], and for the local case at normal surface
singularities see [10, 15].

Algebraic stability is valuable in particular because the first dynamical
degree becomes easy to compute. This can be defined by
\[
 \lambda_1(f)=\lim_{n\to\infty}\|(f^n)^*\|^{1/n}.
\]
For an $f:X\rat X$ which is not algebraically stable, this computation is
usually intractable. However, if one has a stabilisation
$\phi:(g,Y)\rat(f,X)$ then the first dynamical degree, which is birationally
invariant, can be computed simply as the spectral radius of $g^*$. This
also shows that rational maps which can be stabilised have a first dynamical
degree that is an algebraic integer. The recent work of Bell, Diller, and
Jonsson [4] gives an example of a rational map with transcendental first
dynamical degree, hence providing another rational map which could never
be stabilised.

In any case the arguments in Diller--Favre support the idea that blowing
up destabilising orbits is a good approach to achieving algebraic stability.
The evidence of various examples [1, 2, 3, 9] shows that this approach not
only succeeds in many cases but is practical to carry out explicitly.

Let us call a destabilising orbit minimal when it does not contain any
shorter ones.

\sourcestatement{Proposition 1.2}{Suppose that $f:X\rat X$ is a rational
map on a surface. Let $p,f(p),\ldots,f^{n-1}(p)$ be a minimal destabilising
orbit for $f$ and $\pi:X'\to X$ be the birational morphism blowing up each
$p_j=f^{j-1}(p)$. Then any birational morphism $\rho:(g,Y)\to(f,X)$
stabilising $(f,X)$ factors as $\rho=\pi\circ\nu$ for some birational
morphism $\nu:Y\to X'$.}

% SOURCE physical 4, printed 49: complete algorithm and main statement.
\sourceheading{Definition 1.3 (Minimal Stabilisation Algorithm)}
\textit{Given a rational surface map $f_0:X_0\rat X_0$ we define a
(possibly finite) sequence
$\pi_m:(f_{m+1},X_{m+1})\to(f_m,X_m)$ for $m\geq0$ as follows}
\begin{enumerate}
 \item \textit{If $f_m$ is algebraically stable, stop.}
 \item \textit{If not, then pick a minimal destabilising orbit
 $p_1,p_2,\ldots,p_n$ and blowup each of the $p_j$ to produce
 $\pi_m:(f_{m+1},X_{m+1})\to(f_m,X_m)$.}
\end{enumerate}
\textit{If this sequence terminates at $f_M:X_M\rat X_M$, write
$\pi=\pi_1\circ\cdots\circ\pi_{M-1}:X_M\to X$. Then $(f_M,X_M)$ is
algebraically stable and we call $\pi:(f_M,X_M)\to(f,X)$ a minimal
stabilisation of $(f,X)$ (by blowups) when it exists.}

This terminology is justified by the next theorem, which says that the
final result of the Minimal Stabilisation Algorithm, when it terminates,
has a universal property.

\sourcestatement{Theorem 1.4}{Let $f:X\rat X$ be a rational map on a
surface. If there exists a birational morphism $\rho:(g,Y)\to(f,X)$
stabilising $f$ then any instance of the Minimal Stabilisation Algorithm
terminates in a minimal stabilisation
$\pi:(\widehat f,\widehat X)\to(f,X)$ such that $\rho=\pi\circ\nu$ for
some $\nu:Y\to\widehat X$. It follows that the minimal stabilisation
$(\widehat f,\widehat X)$ is unique for $(f,X)$.}

\sourcestatement{Corollary 1.5}{Let $f:X\rat X$ be a birational map on
a surface. Then there exists a unique minimal stabilisation
$\pi:(\widehat f,\widehat X)\to(f,X)$.}
\begin{proof}
By [7], there exists a birational morphism $\rho:\widetilde X\to X$ which
makes $\widetilde f=\rho^{-1}\circ f\circ\rho$ algebraically stable on
$\widetilde X$. By Theorem 1.4, the minimal stabilisation
$\widehat f:\widehat X\rat\widehat X$ via $\pi:\widehat X\to X$ exists
(uniquely and factors $\rho$).
\end{proof}

\section*{Author's text: 2. Background}
% SOURCE physical 7, printed 52: paragraph introducing Section 2.
For the rest of this article, assume all surfaces are smooth projective
over an algebraically closed field, and rational maps are dominant. An
account of most of the facts below can be found in [17, \S V].

% SOURCE physical 8, printed 53: complete notation and graph context.
Let $X,Y$ be surfaces and $f:X\rat Y$ a rational map. Let $U$ be the
largest (open) set on which $f:U\to Y$ is a morphism, then we define the
indeterminate set as $I(f)=X\setminus U$. Alternatively, these are the
finitely many points at which $f$ cannot be continuously defined. These
are often also called fundamental points.

An irreducible curve $C\subset X$ is exceptional iff $f(C\setminus I(f))$
is a point in $Y$. We define the exceptional set, $\mathcal E(f)$, of $f$
to be the union of all (finitely many) irreducible exceptional curves in
$X$. Denote by $\exc(f)$ the number of irreducible components in
$\mathcal E(f)$.

\sourceheading{Remark 2.1}
In the case of the inverse of a birational map, $I(f^{-1})$ is the proper
transform of $\mathcal E(f)$, or the set of values of $f$ with infinite
preimage. Using this latter definition we generalise the meaning of
$I(f^{-1})$ to any rational map $f$.

Let $f:X\rat Y$ be a rational map of surfaces. The graph of $f$ is the
subvariety
\[
 \Gamma_f=\overline{\{(x,f(x))\in X\setminus I(f)\times Y\}}
 \subset X\times Y
\]
along with projections $\alpha:\Gamma_f\to X$ onto the first factor and
$\beta:\Gamma_f\to Y$ onto the second factor which are proper.
$\Gamma_f$ is irreducible because $X\setminus I(f)$ is.

$I(f)$ is the set of points where $\alpha$ does not have a local inverse.
In particular $\alpha^{-1}:X\setminus I(f)\to\Gamma_f$ is an isomorphism.
For $p\notin I(f)$ we have $\beta(\alpha^{-1}(p))=f(p)$. In general for
any set of points $S\subseteq X$, we may define the total transform of $S$
by $f$ as $f(S)=\beta(\alpha^{-1}(S))$. When
$\varnothing\neq S\subseteq I(f)$ this image has dimension 1.
$\mathcal E(f)\subset X$ is the $\alpha$ projection of the set of points
where $\beta$ is not finite-to-one.

$\Gamma_f$ need not be smooth and it will be more convenient to work with
the minimal smooth desingularisation $\Sigma_f$ of $\Gamma_f$. Abusing
notation slightly, we denote the lifted projections as
$\alpha:\Sigma_f\to X$ and $\beta:\Sigma_f\to Y$. Recall the
Néron--Severi group, $\NS(X)$; when $Y=X$ and $f$ is birational, $\beta$
is also a birational morphism, and one can deduce that
$\exc(\alpha)=\exc(\beta)=\operatorname{rk}\NS(\Sigma_f)-\operatorname{rk}\NS(X)$.

\sourceheading{Remark 2.2}
Equivalently $\Sigma_f$ and $\alpha:\Sigma_f\to X$ can be obtained as the
birational morphism, blowing up $X$, which (minimally) resolves the
indeterminacy of $f$. Once resolved $f:X\rat Y$ will lift through
$\alpha$ to a morphism which is precisely $\beta:\Sigma_f\to Y$.
We obtain a map $\alpha\times\beta:\Sigma_f\to X\times Y$ and one can
check its image is $\Gamma_f$, hence it is a smooth desingularisation which
dominates the minimal one. Since the minimal desingularisation is a
candidate for resolving the indeterminacy we also get a birational morphism
the other direction, meaning the two definitions must agree.

% SOURCE physical 9, printed 54.
\sourceheading{Remark 2.3 (Warning)}
$\Sigma_f$ may contain curves which appear neither in the domain $X$ or
the codomain $Y$, meaning both $\alpha$ and $\beta$ map the curve to a
point. This issue will be rectified in Proposition 4.2 and Corollary 4.6.
For an example, pick any surface with a rational curve $C$ of
self-intersection 0, and construct the birational map which blows up a
point on $C$ to give an exceptional curve $D$, blows up another point of
$D$ to give $E$, but then blows down $C$ followed by $D$. The second curve
$D$ (i.e. its proper transform) is contracted in both the domain and range
of this map, however it can be found in the smooth graph (which is the
intermediate surface containing all three curves).

\sourcestatement{Proposition 2.4 (See [17, \S V 5.3 \& 5.4])}{Let
$f:X\to Y$ be a birational morphism of surfaces. Then $f$ can be written
as a composition of $\exc(f)$ point blowups.}
\textit{Suppose $p\in I(f^{-1})$ and $\pi:Y'\to Y$ be the point blowup of
$p$. Then $f$ factors as $\pi\circ f'$ where $f':X\to Y'$ is a birational
morphism with $\exc(f')=\exc(f)-1$.}

\textit{Otherwise if $p\notin I(f^{-1})$ and $\rho:X'\to X$ is the point
blowup of $f^{-1}(p)$. Then $f$ lifts to a birational morphism
$f'=\pi^{-1}\circ f\circ\rho$ with $\exc(f')=\exc(f)$.}

\sourceheading{Remark 2.5}
Note that Proposition 2.4 does not hold for all rational morphisms, meaning
it can fail if we replace ``birational'' with ``rational'', and consider
$I(f^{-1})$ as defined in Remark 2.1. The second part of the proposition
fails in Example 6.9.

\sourcestatement{Definition 2.6}{Let $f:X\rat Y$ and $g:Y\rat Z$ be
rational maps on surfaces. We say an irreducible curve $C\subset X$ is a
destabilising curve for the composition $g\circ f$ iff
$f(C\setminus I(f))=y\in I(g)$. Equivalently $C\subseteq f^{-1}(y)$ and
$g(y)\supseteq D$ for some irreducible curve $D\subset Z$. We call $D$ an
inverse destabilising curve, and we say that $y$ destabilises the composition
$g\circ f$. We say that the composition $g\circ f$ is locally stable at
$x\in X$ iff $x$ is not contained in any destabilising curve.}

By unravelling the definitions given above, we get the following proposition.

\sourcestatement{Proposition 2.7}{Let $f:X\rat Y$ and $g:Y\rat Z$ be
rational maps. Then
\[
 g(f(x))\supseteq(g\circ f)(x)
\]
with equality if $x$ is not contained in a destabilising curve. Moreover
$z\in g(f(x))$ if and only if $z\in(g\circ f)(x)$ or we can find
$y\in I(g)\cap I(f^{-1})$ destabilising the composition $g\circ f$ where
$z\in g(y)$ and $x\in f^{-1}(y)$.}

% SOURCE physical 10, printed 55: complete definitions and criterion.
\sourcestatement{Definition 2.8}{Let $f:X\rat X$ be a rational map and
let $p,f(p),\ldots,f^{m-1}(p)$ be an orbit in $X$. We say this is a
destabilising orbit iff there is an (indeterminate) $q\in f^{m-1}(p)$ such
that both $f(q)=D$ and $f^{-1}(p)=C$ are curves. Let $N\in\mathbb N$, then
we say the orbit is $N$-eventually destabilising iff there is a closed point
$q\in f^{m-1}(p)$ such that both $f(q)=D$ and $f^{-n}(p)=C$ are curves
for some $n\geq N$.}
\textit{In either destabilising scenario, we say the length of the orbit is
$m$, we call each irreducible component of $C$ a destabilising curve and
each component of $D$ an inverse destabilising curve of $f$.}

\sourcestatement{Definition 2.9}{Let $f:X\rat X$ be a rational map. We
say that $f$ is algebraically stable iff for every $n\geq0$ we have
$(f^n)^*=(f^*)^n$. We say $f$ is $N$-eventually algebraically stable iff
for every $n\geq N$ and $m\geq0$ we have
$(f^{n+m})^*=(f^n)^*(f^m)^*=(f^n)^*(f^*)^m$.}

\sourcestatement{Proposition 2.10 ([7, 1.14], [13])}{Let $f:X\rat X$ be
a rational map. Then $f$ is algebraically stable if and only if $f$ has
no destabilising orbits. Similarly, $f$ is $N$-eventually algebraically
stable if and only if $f$ has no $N$-eventually destabilising orbits.}

This geometric characterisation of algebraic stability can be understood
through Proposition 2.7. One can show that $f^*g^*D-(g\circ f)^*D$ is
supported on destabilising curves for any divisor $D$ and moreover that
the composition $g\circ f$ is stable if and only if $f^*g^*=(g\circ f)^*$.

\section*{Author's text: 3. Minimal Stabilisation}
\sourcestatement{Definition 3.1}{We say that a destabilising orbit
$p,f(p),\ldots,f^{m-1}(p)$ is minimal iff the $p_j=f^{j-1}(p)$ are all
distinct closed points where for $1\leq j<m$ we have $p_j\notin I(f)$,
and for $1<j\leq m$ we have $p_j\notin I(f^{-1})$.}

It is easy to see that the points of a minimal destabilising orbit do not
contain any shorter destabilising orbits and conversely that any
destabilising orbit of minimum length $m$ is minimal, so these must always
exist when $f$ is not algebraically stable. Minimality is less natural for
eventually destabilising orbits. An $N$-eventually destabilising orbit always
possesses a point that constitutes a singleton $N$-eventually destabilising
orbit, which is vacuously minimal.

\begin{proof}[Proof of Proposition 1.2]
First note that by applying Proposition 2.4 ($n$ times), if $\rho$ blows up
all the $p_j$ at least once, then we get a new birational morphism $\nu$
which provides the factorisation $\rho=\pi\circ\nu$. Therefore we will
% SOURCE physical 11, printed 56: proof continues without omission.
proceed to show that $\rho$ does indeed blowup the $p_j$. Let
$\widehat p_j=\rho^{-1}(p_j)$, which a priori may be closed points or
exceptional curves.
\[
\begin{tikzcd}[row sep=large,column sep=large]
Y \arrow[r,dashed,"g"] \arrow[d,"\rho"'] & Y \arrow[d,"\rho"] \\
X \arrow[r,dashed,"f"] & X
\end{tikzcd}
\]
Suppose not; then there is a largest $m\leq n$ such that $\widehat p_m$
is a closed point in $Y$, so we claim this is indeterminate for $g$. Say
$m=n$ and $f(p_n)=D\subset X$, then $\widehat p_n\in I(g)$ if
$g(\widehat p_n)$ is a curve, for which it is enough to show that
$\rho(g(\widehat p_n))$ is a curve. Indeed
\[
\rho(g(\widehat p_n))=\rho\circ g(\widehat p_n)
=f\circ\rho(\widehat p_n)=f(p_n)=D
\]
is a curve; note that $I(\rho)=\varnothing$ and the last step holds because
$\rho$ is locally an isomorphism at $\widehat p_n$. If $m<n$ then
$p_{m+1}\notin I(f^{-1})$, so the composition $\rho^{-1}\circ f$ is locally
stable at $p_m$, and $\rho$ is locally an isomorphism at $\widehat p_m$
so $(\rho^{-1}\circ f)\circ\rho=g$ is locally stable near
$\widehat p_m$. Therefore
\[
g(\widehat p_m)=\rho^{-1}(f(\rho(\widehat p_m)))=\rho^{-1}(p_{m+1})
\]
which is also a curve by assumption.

Suppose that $k\leq m$ is minimal such that $p_k,\ldots,p_m$ are not
blown up by $\rho$ and either $k=1$ or $k-1$ is blown up by $\rho$. Since
$\rho$ is a local isomorphism over these points, one sees that
$\widehat p_k\mapsto\widehat p_{k+1}\mapsto\cdots\mapsto\widehat p_m$.
To provide a contradiction to the algebraic stability of $g$, we will next
show that this is a destabilising orbit for $g$. We only need to show
that $g^{-1}(p_k)$ is a curve and this is very similar to the case of $p_m$
above.

$\rho$ is locally an isomorphism near $\widehat p_k$. So it is enough to
show that $(f\circ\rho)^{-1}(p_k)$ is a curve. Say $k=1$ and hence
$f^{-1}(p_1)$ contains a curve $C$, then $(f\circ\rho)^{-1}(p_1)$ contains
$\rho^{-1}(C\setminus I(f))$ which is (almost all of) a curve. Otherwise
$k>1$, so $p_{k-1}\notin I(f)$ and the composition $f\circ\rho$ is
algebraically stable over $\rho^{-1}(p_{k-1})$ which is a curve
$\widehat C$. We get that $\rho^{-1}\circ f^{-1}(p_k)=\widehat C$.
\end{proof}

\begin{proof}[Proof of Theorem 1.4]
Given that $g:Y\rat Y$ dominates $f$ via
$\rho:(g,Y)\to(f,X)=(f_1,X_1)$ we may proceed inductively on $m$ with the
hypothesis that $g:Y\rat Y$ dominates $f_m:X_m\rat X_m$ via
$\nu_m:Y\to X_m$.

If $f_m$ is algebraically stable we are done, otherwise Proposition 1.2
says that because $\pi_m$ blows up a minimal destabilising orbit we have
a $\nu_{m+1}:Y\to X_{m+1}$ which factors $\nu_m$ as
$\nu_m=\pi_m\circ\nu_{m+1}$.
% SOURCE physical 12, printed 57: final proof, all paragraphs and diagram.
\[
\begin{tikzcd}[row sep=large,column sep=large]
X_{m+1} \arrow[d,"\pi_m"'] & Y \arrow[l,red,"\nu_{m+1}"'] \arrow[dl,"\nu_m"] \\
X_m & {}
\end{tikzcd}
\]
Clearly there is no limit to the number of times we can do this if $f_m$
is never algebraically stable for $m\geq1$. However overall we have shown
that
\[
\rho=\nu_1=\pi_1\circ\nu_2=\pi_1\circ\pi_2\circ\nu_3
=\cdots=\pi_1\circ\cdots\circ\pi_m\circ\nu_{m+1}=\pi\circ\nu_{m+1}
\]
meaning that $\exc(\rho)\geq m$. Therefore
$\{m\in\mathbb N:f_m\text{ is not AS}\}$ is in fact bounded above,
strictly by $m=M$ say, and whence $f_M$ is algebraically stable. Moreover,
$\rho=\pi\circ\nu$ where we define $\nu=\nu_M$.

We have shown that $\pi$ factors any birational morphism stabilising $f$,
and to finish we apply this to get uniqueness of $(\widehat f,\widehat X)$.
Suppose we proceed in the Minimal Stabilisation Algorithm in two different
ways which produce two (potentially different) models, namely
$\widehat f_1:\widehat X_1\rat\widehat X_1$ via
$\pi_1:\widehat X_1\to X$ and
$\widehat f_2:\widehat X_2\rat\widehat X_2$ via
$\pi_2:\widehat X_2\to X$. By the above we have that
$\pi_1=\pi_2\circ\nu_1$ and $\pi_2=\pi_1\circ\nu_2$. We deduce that
$\nu_1,\nu_2$ are inverse morphisms to each other, providing an
isomorphism not only of surfaces but dynamical systems
$\nu_1:(\widehat f_1,\widehat X_1)\leftrightarrow
(\widehat f_2,\widehat X_2)$, i.e. the following diagram commutes.
\[
\begin{tikzcd}[row sep=large,column sep=huge]
\widehat X_1
 \arrow[rr,bend right=15,"\nu_1"']
 \arrow[dr,"\pi_1"']
 \arrow[loop left,"\mathrm{id}"']
 & {} & \widehat X_2
 \arrow[ll,bend right=15,"\nu_2"']
 \arrow[dl,"\pi_2"]
 \arrow[loop right,"\mathrm{id}"] \\
 & X & {}
\end{tikzcd}
\]
\end{proof}

\clearpage
% SOURCE physical 29, printed 74: full bibliography, original numbers.
\begin{thebibliography}{22}
\bibitem{b1} E. Bedford \& K. Kim, ``Periodicities in linear fractional
recurrences: degree growth of birational surface maps'', \emph{Mich. Math. J.}
54 (2006), no. 3, p. 647--670.
\bibitem{b2} E. Bedford \& K. Kim, ``Continuous families of rational surface
automorphisms with positive entropy'', \emph{Math. Ann.} 348 (2010), no. 3,
p. 667--688.
\bibitem{b3} E. Bedford, K. Kim, T. T. Truong, N. Abarenkova \&
J.-M. Maillard, ``Degree complexity of a family of birational maps'',
\emph{Math. Phys. Anal. Geom.} 11 (2008), no. 1, p. 53--71.
\bibitem{b4} J. P. Bell, J. Diller \& M. Jonsson, ``A transcendental dynamical
degree'', \emph{Acta Math.} 225 (2020), no. 2, p. 193--225.
\bibitem{b5} R. L. Benedetto, \emph{Dynamics in one non-archimedean variable},
Graduate Studies in Mathematics, vol. 198, American Mathematical Society,
2019, xviii+463 pages.
\bibitem{b6} R. A. P. Birkett, ``Algebraic Stability and Skew Products on the
Berkovich Projective Line'', PhD Thesis, University of Notre Dame, 2023.
\bibitem{b7} J. Diller \& C. Favre, ``Dynamics of Bimeromorphic Maps of
Surfaces'', \emph{Am. J. Math.} 123 (2001), p. 1135--1169.
\bibitem{b8} J. Diller \& J.-L. Lin, ``Rational surface maps with invariant
meromorphic two-forms'', \emph{Math. Ann.} 364 (2016), no. 1--2, p. 313--352.
\bibitem{b9} C. Favre, ``Les applications monomiales en deux dimensions'',
\emph{Mich. Math. J.} 51 (2003), no. 3, p. 467--475.
\bibitem{b10} C. Favre, ``Holomorphic self-maps of singular rational
surfaces'', \emph{Publ. Mat., Barc.} 54 (2010), no. 2, p. 389--432.
\bibitem{b11} C. Favre \& M. Jonsson, \emph{The valuative tree}, Lecture
Notes in Mathematics, vol. 1853, Springer, 2004, xiv+234 pages.
\bibitem{b12} C. Favre \& M. Jonsson, ``Dynamical compactifications of
$\mathbb C^2$'', \emph{Ann. Math. (2)} 173 (2011), no. 1, p. 211--249.
\bibitem{b13} J. E. Fornæss \& N. Sibony, ``Complex dynamics in higher
dimensions'', in \emph{Complex potential theory (Montreal, PQ, 1993)}, NATO
Adv. Sci. Inst. Ser. C: Math. Phys. Sci., vol. 439, Kluwer Academic Publishers,
1994, p. 131--186.
\bibitem{b14} W. Gignac \& M. Ruggiero, ``Growth of attraction rates for
iterates of a superattracting germ in dimension two'',
\emph{Indiana Univ. Math. J.} 63 (2014), no. 4, p. 1195--1234.
\bibitem{b15} W. Gignac \& M. Ruggiero, ``Local dynamics of non-invertible
maps near normal surface singularities'', \emph{Mem. Am. Math. Soc.} 272
(2021), no. 1337, p. xi+100.
\bibitem{b16} B. Grammaticos, A. Ramani \& V. Papageorgiou, ``Do integrable
mappings have the Painlevé property?'', \emph{Phys. Rev. Lett.} 67 (1991),
p. 1825--1828.
\bibitem{b17} R. Hartshorne, \emph{Algebraic Geometry}, Springer, 1977.
\bibitem{b18} J. H. Hubbard \& P. Papadopol, ``Superattractive fixed points
in $\mathbb C^n$'', \emph{Indiana Univ. Math. J.} 43 (1994), no. 1, p. 321--365.
\bibitem{b19} M. Jonsson, ``Dynamics of Berkovich spaces in low dimensions'',
in \emph{Berkovich spaces and applications}, Lecture Notes in Mathematics,
vol. 2119, Springer, 2015, p. 205--366.
\bibitem{b20} M. Jonsson \& E. Wulcan, ``Stabilization of monomial maps'',
\emph{Mich. Math. J.} 60 (2011), no. 3, p. 629--660.
\bibitem{b21} J.-L. Lin, ``Algebraic stability and degree growth of monomial
maps'', \emph{Math. Z.} 271 (2012), no. 1--2, p. 293--311.
\bibitem{b22} A. Lonjou \& C. Urech, ``Actions of Cremona groups on CAT(0)
cube complexes'', to appear in \emph{Duke Math. J.}, 2021,
\url{https://arxiv.org/abs/2001.00783}.
\end{thebibliography}
\end{document}
